\scanpage{281}
\mspage{197}
% archive sheet number: 280
qui est en effet \unsure{vérifiée|vérifié} par définition de $u$. Cela \unsure{achève|qchève} de prouver $2^\circ$.

Prouvons $3^\circ$. Comme $i'$ est pl. fidèle, en vertu du $1^\circ$, pour prouver que $i'$ est str. générateur, il suffit d'établir que \inserted{pour \unsure{tt} objet $X$ de $M$,} la famille des flèches $i'(x')\longrightarrow X$, indexée par $\mathrm{Ob}\,C'_{/X}$, est épimorphique stricte.

Or cette famille contient la famille des $\underbrace{i'(\alpha(x))}_{\substack{\Vert\\i(x)}}\longrightarrow X$, \deleted{\unsure{$x$}} \inserted{indexée par $\mathrm{Ob}\,C_{/X}$.} \deleted{\unsure{cette dernière étant épi}} Donc il résulte de l'hypothèse sur $i$, et du $2^\circ$, que cette famille est épimorphique stricte. On a gagné !

Ouf\unsure{...|.--} \par\hfill\unsure{$\backslash$}\hfill\null

\uline{Démonstration de 4.8.1}. On a

\[ c'\Longrightarrow c\Longrightarrow b\Longrightarrow a, \]

la première implication \unsure{tautologique}, la seconde résulte du fait que $\unsure{\pi|\kappa|k}_0$ \unsure{cardinal|card.} régulier pour la fega canonique, implique que

\[ M\xleftarrow{\approx}\mathrm{Ind}_{\pi_0}(\unsure{F|\mathcal F}^{\pi_0}), \]

$F^{\pi_0}$ équivalente à une petite catégorie \inserted{$C$}, d'où b). \deleted{\unsure{ou}}
